\documentclass[11pt,oneside]{book}

% ---------------------------------------------------------------------------
% Fonts and language
% ---------------------------------------------------------------------------
\usepackage[T1]{fontenc}
\usepackage[utf8]{inputenc}
\usepackage[main=brazil]{babel}
% tcolorbox (and its pgf/tikz dependency) must load early, right after
% babel -- loading it after several other packages under the `book` class
% exhausts a TeX register pool that pgfcorearrows.code.tex needs, producing
% a "Missing \endcsname inserted" error deep inside pgf itself.
\usepackage[most]{tcolorbox}
\usepackage[strings]{underscore}
\usepackage{csquotes}
\usepackage{tgpagella}      % body serif -- a Palatino-like humanist face
\usepackage{tgheros}        % sans, used only for numerals and labels
\usepackage{tgcursor}       % mono, used only for code identifiers
\usepackage{amsmath}
\usepackage{microtype}
\frenchspacing

% ---------------------------------------------------------------------------
% Page geometry
% ---------------------------------------------------------------------------
\usepackage{geometry}
\geometry{a4paper,margin=2.8cm,bindingoffset=0mm}

% ---------------------------------------------------------------------------
% Colour palette -- warm, quiet, matches the instrument's own wood tones
% ---------------------------------------------------------------------------
\usepackage{xcolor}
\definecolor{ink}{HTML}{1B1B1B}
\definecolor{accent}{HTML}{9C5330}
\definecolor{rulecol}{HTML}{D8CFC4}
\definecolor{soft}{HTML}{6B6B6B}

% ---------------------------------------------------------------------------
% Lists, tables, boxes, drop caps
% ---------------------------------------------------------------------------
\usepackage{enumitem}
\setlist{itemsep=3pt,topsep=5pt,parsep=0pt}
\renewcommand\labelitemi{\color{accent}\textbullet}
\usepackage{booktabs}
\usepackage{tabularx}
\usepackage{lettrine}
\setlength{\DefaultFindent}{2pt}

% ---------------------------------------------------------------------------
% Headers, footers
% ---------------------------------------------------------------------------
\usepackage{fancyhdr}
\pagestyle{fancy}
\fancyhf{}
\fancyhead[C]{\small\sffamily\color{soft}\nouppercase{\leftmark}}
\fancyfoot[C]{\small\color{soft}\thepage}
\renewcommand{\headrulewidth}{0pt}
\renewcommand{\chaptermark}[1]{\markboth{#1}{}}
\fancypagestyle{plain}{%
  \fancyhf{}
  \fancyfoot[C]{\small\color{soft}\thepage}
  \renewcommand{\headrulewidth}{0pt}
}

% ---------------------------------------------------------------------------
% Chapter and section titles -- quiet numeral, then the title, then a
% hairline rule. Nothing else fights for attention on the page.
% ---------------------------------------------------------------------------
\usepackage{titlesec}

\titleformat{\chapter}[display]
  {\normalfont}
  {\sffamily\color{accent}\Large\thechapter}
  {6pt}
  {\raggedright\huge\bfseries\color{ink}}
  [\vspace{6pt}{\color{rulecol}\titlerule}\vspace{20pt}]
\titlespacing*{\chapter}{0pt}{0pt}{10pt}

\titleformat{name=\chapter,numberless}[display]
  {\normalfont}
  {}
  {0pt}
  {\raggedright\huge\bfseries\color{ink}}
  [\vspace{6pt}{\color{rulecol}\titlerule}\vspace{20pt}]

\titleformat{\section}
  {\normalfont\sffamily\bfseries\color{ink}\large}{\thesection}{8pt}{}
\titlespacing*{\section}{0pt}{22pt}{8pt}

\titleformat{\subsection}
  {\normalfont\sffamily\bfseries\color{ink}}{}{0pt}{}
\titlespacing*{\subsection}{0pt}{16pt}{6pt}

% ---------------------------------------------------------------------------
% Links
% ---------------------------------------------------------------------------
\usepackage[colorlinks=true,linkcolor=accent,urlcolor=accent,citecolor=accent,
  bookmarksopen=true,pdfstartview=Fit]{hyperref}

% ---------------------------------------------------------------------------
% The "new word" box -- a single quiet visual device, reused for every
% definition and every honesty note, deliberately not decorated further.
% ---------------------------------------------------------------------------
\newtcolorbox{newword}{
  colback=white, colframe=accent, boxrule=0pt, leftrule=1.1pt,
  arc=0pt, outer arc=0pt,
  left=12pt, right=8pt, top=7pt, bottom=7pt,
  before skip=12pt, after skip=12pt,
  fontupper=\small
}
\newcommand{\nova}[1]{\textsc{\color{accent}#1.}~}

% Drop-cap helper, used once per chapter opening
\newcommand{\opening}[2]{\lettrine[lines=2,loversize=0.08,lraise=0.02]{#1}{#2}}

\newcommand{\crate}[1]{\texttt{\small #1}}

\title{Como este piano funciona}
\author{}
\date{}

\begin{document}

% =============================================================================
\frontmatter
\pagestyle{empty}

% -- Title page --------------------------------------------------------------
\begin{titlepage}
\centering
\vspace*{5cm}
{\sffamily\color{soft}\small\MakeUppercase{Uma aula completa, sem pré-requisitos}}\\[1.2cm]
{\sffamily\bfseries\color{ink}\Huge Como este piano\\[2pt]funciona}\\[0.9cm]
{\color{rulecol}\rule{5cm}{0.6pt}}\\[0.9cm]
{\itshape\color{soft}\large A física e a matemática por trás de cada nota,\\
explicadas do zero, em português.}
\vfill
{\sffamily\small\color{soft}Desenvolvimento patrocinado pela \href{https://grabatus.com}{Grabatus Labs}}\\[4pt]
{\sffamily\footnotesize\color{soft}\MakeUppercase{piano --- um sintetizador de piano fisicamente modelado, em Rust}}
\vspace*{1.5cm}
\end{titlepage}

% -- Colophon ------------------------------------------------------------
\thispagestyle{empty}
\vspace*{2cm}
\noindent{\small\color{soft}
\textsc{Como este piano funciona} é a versão didática, sem pré-requisitos,
da documentação técnica do projeto \emph{piano} --- um sintetizador de piano
fisicamente modelado, escrito em Rust, sem nenhuma amostra ou gravação de
áudio real em nenhum lugar do código.

\vspace{8pt}
O código-fonte, a documentação técnica completa (em inglês) e este mesmo
documento em Markdown estão publicados abertamente em
\url{https://github.com/rodolphomacedo/piano}, sob licença MIT ou
Apache-2.0, à escolha de quem usar.

\vspace{8pt}
Desenvolvimento patrocinado pela Grabatus Labs --- \url{https://grabatus.com}.

\vspace{8pt}
Tipografado com \LaTeX{} (motor Tectonic), nas famílias TeX Gyre Pagella,
TeX Gyre Heros e TeX Gyre Cursor.
}

% -- Before we begin ----------------------------------------------------
\chapter*{Antes de começar}
\addcontentsline{toc}{chapter}{Antes de começar}
\thispagestyle{plain}

\noindent\textbf{Para quem é este documento?} Para qualquer pessoa ---
inclusive quem nunca estudou física, música ou programação. Vamos construir
cada ideia do zero, devagar, com exemplos do dia a dia antes de qualquer
fórmula. Se em algum momento aparecer uma palavra nova (de música, física ou
matemática), ela é explicada \emph{antes} de ser usada, com uma caixa como
esta:

\begin{newword}
\nova{Palavra nova} Explicação simples, com exemplo.
\end{newword}

\noindent Não existe pré-requisito. Se você conseguiu ler até aqui, consegue
ler o documento inteiro.

\tableofcontents

% =============================================================================
\mainmatter
\pagestyle{fancy}

% =============================================================================
\chapter{O que este projeto realmente faz}

\opening{E}{xistem} dois jeitos completamente diferentes de fazer um
\enquote{piano digital} tocar uma nota.

\textbf{Jeito 1 --- gravar e reproduzir.} Alguém senta em um piano de
verdade, grava cada tecla sendo tocada em vários volumes diferentes, e
guarda milhares de arquivos de áudio. Quando você aperta uma tecla no
teclado digital, ele simplesmente \emph{toca de volta} a gravação certa. É
assim que a imensa maioria dos pianos digitais, teclados e aplicativos de
celular funciona --- isso se chama \emph{sampler}.

\begin{newword}
\nova{Sampler} Um instrumento que reproduz gravações prontas em vez de
calcular o som.
\end{newword}

Funciona bem, mas tem um problema conceitual: o piano nunca está
\enquote{tocando} de verdade, ele está apenas apertando o \enquote{play} numa
gravação antiga.

\textbf{Jeito 2 --- calcular a física.} É o que este projeto faz. Dentro do
código não existe nenhum arquivo de áudio, nenhuma gravação, nenhum
\enquote{play}. Existe uma simulação matemática de como uma corda de piano de
verdade se comporta: como ela vibra quando é golpeada, como perde energia
com o tempo, como duas ou três cordas da mesma tecla interagem, como a caixa
de madeira do piano colore o som. Cada vez que você aperta uma tecla, o
programa \textbf{calcula}, do zero, amostra por amostra, 48 mil vezes a cada
segundo, o que uma corda de piano de verdade faria fisicamente. Isso se
chama \emph{modelagem física}.

\begin{newword}
\nova{Modelagem física} (\emph{physical modelling}, em inglês) Descrever um
objeto do mundo real --- aqui, uma corda de piano --- como um conjunto de
equações, e resolver essas equações em tempo real para produzir o
comportamento do objeto, em vez de gravar o objeto.
\end{newword}

Essa é a razão de existir de todo este documento: explicar, devagar e com
exemplos, \textbf{como é possível calcular o som de uma corda de piano em
tempo real}, e \textbf{quais pedacinhos da física real} este projeto já
modelou (e quais ainda faltam).

Nada aqui é segredo ou \enquote{mágica de IA}. É física do século XIX e XX
(a mesma que rege cordas de violão, o balanço de um parquinho e as ondas do
mar), descrita como fórmulas, e as fórmulas transformadas em código Rust que
roda rápido o suficiente para acontecer ao vivo.

% =============================================================================
\chapter{O que é som, afinal}

\opening{A}{ntes} de falar de piano, precisamos entender o que é som ---
porque tudo o que vem depois é, no fundo, sobre como \emph{fabricar} som a
partir de números.

Bata palmas. O que aconteceu, fisicamente? Suas mãos empurraram o ar que
estava entre elas. Esse ar empurrado empurrou o ar vizinho, que empurrou o
próximo, e assim por diante --- uma onda de pressão se espalhando pelo
ambiente, exatamente como quando você joga uma pedra num lago parado e vê os
círculos se espalharem na água.

\begin{newword}
\nova{Onda} Uma perturbação que se propaga através de um meio (ar, água,
uma corda) sem que o próprio meio \enquote{viaje} --- é a \emph{perturbação}
que anda, não o ar em si. Pense em uma ola de estádio: cada torcedor só
levanta e senta no próprio lugar, mas o movimento parece \enquote{correr}
pela arquibancada.
\end{newword}

Quando essa onda de pressão chega até o seu ouvido, ela empurra e puxa o
tímpano --- uma membrana bem fininha, como a pele esticada de um tambor ---
para frente e para trás, muito rapidamente. O cérebro interpreta esse
vaivém-vaivém-vaivém como som.

A pergunta importante é: \textbf{com que velocidade} o ar empurra e puxa? Se
você bater palmas devagar (uma vez por segundo), o ar até vibra, mas tão
devagar que você não ouviria \enquote{som} nenhum --- ouviria só um
\enquote{bump} isolado. Para existir uma nota musical, o vaivém precisa se
repetir \emph{centenas ou milhares de vezes por segundo}. É disso que o
próximo capítulo trata.

% =============================================================================
\chapter{O que é uma nota musical}

\begin{newword}
\nova{Frequência} Quantas vezes por segundo algo se repete. Se seu coração
bate 70 vezes por minuto, a frequência dos seus batimentos é \enquote{70 por
minuto}. Para som, medimos frequência em \textbf{hertz} (abreviado
\textbf{Hz}), que quer dizer \enquote{vezes por segundo}. Um som de 440~Hz
vibra o ar para frente e para trás 440 vezes a cada segundo.
\end{newword}

\opening{U}{ma} \textbf{nota musical} nada mais é do que um som cuja
frequência é constante e bem definida --- o vaivém do ar se repete no mesmo
ritmo, sem parar, durante o tempo em que a nota soa. Quanto \textbf{mais
alta} a frequência, mais \textbf{aguda} (fina) a nota soa aos nossos
ouvidos; quanto \textbf{mais baixa}, mais \textbf{grave} (grossa).

Um piano tem 88 teclas, e cada uma corresponde a uma frequência diferente. A
tecla mais grave (a mais à esquerda, chamada \textbf{A0}, ou \textbf{Lá0})
vibra a 27,5~Hz. A tecla mais aguda (a mais à direita, \textbf{C8}, ou
\textbf{Dó8}) vibra a 4186~Hz --- mais de 150 vezes mais rápido. A nota de
referência que afinadores de piano (e este projeto) usam como ponto de
partida é o \textbf{A4} (\textbf{Lá4}, o Lá \enquote{do meio}), fixado
internacionalmente em \textbf{440~Hz}.

\begin{newword}
\nova{Oitava} Duas notas estão \enquote{uma oitava} de distância quando a
frequência de uma é exatamente o dobro da outra. A4 é 440~Hz; uma oitava
acima, A5, é 880~Hz; uma oitava abaixo, A3, é 220~Hz. É por isso que dobrar
(ou cortar pela metade) a frequência soa, aos nossos ouvidos, como \enquote{a
mesma nota, só que mais aguda/grave} --- não é coincidência, é como o
ouvido humano processa proporções de frequência.
\end{newword}

\begin{newword}
\nova{Notação de nota + oitava (ex.\ A4, C3, F\#5)} A letra é o nome da nota
(A, B, C, D, E, F, G --- o equivalente em inglês de Lá, Si, Dó, Ré, Mi, Fá,
Sol) e o número indica em qual oitava ela está --- quanto maior o número,
mais aguda. Um \texttt{\#} (sustenido) depois da letra significa \enquote{meio
tom acima}. Este projeto (e a maior parte do software musical no mundo) usa
essa notação em inglês, então \enquote{A4} é o Lá central e \enquote{C4} é o
Dó central.
\end{newword}

O piano é um instrumento \textbf{temperado}: o espaço entre A0 e C8 é
dividido em 88 degraus (chamados \textbf{semitons}) espaçados de um jeito
matematicamente preciso, onde cada semitom multiplica a frequência anterior
por aproximadamente $1{,}0595$ (a raiz 12\textsuperscript{a} de 2 --- porque
12 semitons formam uma oitava, e uma oitava dobra a frequência:
$1{,}0595^{12} \approx 2$). Você não precisa decorar esse número; o que
importa entender é que \textbf{cada tecla do piano é uma frequência-alvo
específica e fixa}, e é exatamente essa frequência-alvo que o programa
recebe quando você aperta uma tecla --- o trabalho de todo o resto deste
documento é: \emph{dada essa frequência, como calcular o som de uma corda
vibrando nela?}

% =============================================================================
\chapter{Como um piano de verdade produz som}

\opening{A}{bra} a tampa de um piano de verdade (ou imagine um cortado ao
meio) e você vai ver, para a maioria das teclas, isto:

\begin{enumerate}[leftmargin=1.6em]
\item Você aperta uma tecla.
\item Um mecanismo mecânico (a \textbf{ação}) lança um pequeno
\textbf{martelo coberto de feltro} contra uma ou mais \textbf{cordas de
aço} esticadas sob tensão.
\item O martelo bate na corda e volta --- ele não fica encostado, é uma
martelada rápida, de menos de 5 milésimos de segundo.
\item A corda, agora vibrando, empurra a \textbf{caixa de ressonância} (o
grande tampo de madeira embaixo/atrás das cordas) através de uma peça de
madeira chamada \textbf{ponte}.
\item A caixa de ressonância, por ser grande e leve, empurra o ar de
verdade e é o que você realmente ouve --- a corda sozinha, fininha como é,
mal move ar suficiente para ser audível.
\item Enquanto a tecla continua pressionada, a corda continua vibrando e
perdendo energia aos poucos, até o som sumir. Quando você solta a tecla, um
\textbf{abafador} (\emph{damper}) de feltro encosta na corda e para a
vibração imediatamente.
\item Se você pisar no \textbf{pedal direito} (pedal de sustentação,
também chamado de \emph{sustain}), todos os abafadores do piano levantam
ao mesmo tempo --- a corda que você tocou continua soando mesmo depois de
soltar a tecla, \emph{e} todas as outras cordas do piano ficam livres para
vibrar também, se algo as excitar (capítulo~13).
\end{enumerate}

\begin{newword}
\nova{Ponte} (\emph{bridge}) A peça que transmite a vibração da corda para
a caixa de ressonância, como o \enquote{sela} de um violão.
\end{newword}

Esse é o piano inteiro, em sete passos. Cada capítulo a partir daqui pega
um desses passos e explica: (a) a física real por trás dele, e (b) como
este projeto o transforma em cálculo matemático que um computador consegue
fazer 48 mil vezes por segundo.

% =============================================================================
\chapter{Por que instrumentos diferentes soam diferentes: harmônicos e timbre}

\opening{T}{oque} a mesma nota --- digamos, A4 (440~Hz) --- num piano e num
violino. Ambos vibram a 440~Hz. Ambos são, estritamente, \enquote{a mesma
nota}. E ainda assim eles soam completamente diferentes. Por quê?

A resposta é que \textbf{nenhuma corda vibra em apenas uma frequência}.
Quando uma corda é golpeada ou puxada, ela vibra simultaneamente:

\begin{itemize}
\item inteira, de ponta a ponta (essa é a \textbf{frequência fundamental}
--- a nota \enquote{principal} que você identifica, 440~Hz no nosso
exemplo);
\item em duas metades, cada metade vibrando ao dobro da velocidade
(880~Hz);
\item em três terços, cada terço vibrando ao triplo da velocidade
(1320~Hz);
\item e assim por diante, em infinitas divisões cada vez mais fracas.
\end{itemize}

\begin{newword}
\nova{Harmônico} (ou \emph{parcial}, ou \emph{overtone}) Cada uma dessas
vibrações simultâneas de uma corda. O primeiro harmônico (a corda inteira)
é a fundamental; o segundo, terceiro, quarto harmônico etc. são múltiplos
inteiros dela ($2\times$, $3\times$, $4\times$ a frequência fundamental) e
em geral soam mais fracos quanto mais alto o número.
\end{newword}

\begin{newword}
\nova{Timbre} A \enquote{cor} ou \enquote{textura} característica de um
instrumento --- aquilo que faz um piano soar diferente de um violino
tocando a mesma nota. O timbre é determinado principalmente pela
\textbf{proporção de volume entre os harmônicos}: um piano tem harmônicos
numa certa proporção, um violino tem em outra, e é essa \enquote{receita}
de harmônicos que seu ouvido reconhece como \enquote{isto é um piano}.
\end{newword}

Todo o trabalho de simular um instrumento de verdade se resume, em grande
parte, a produzir a \emph{mistura certa de harmônicos, no volume certo,
morrendo na velocidade certa} --- porque é essa mistura que faz a diferença
entre \enquote{um bipe eletrônico monótono} e \enquote{aquilo soa como um
piano de verdade}. Os capítulos~9 a~14 mostram, um a um, os mecanismos
físicos reais que moldam essa mistura num piano de verdade --- e como este
projeto os reproduz.

% =============================================================================
\chapter{Como um computador \enquote{ouve} e \enquote{fala} som}

\opening{U}{m} computador não entende \enquote{ondas contínuas de pressão
do ar} --- ele só sabe guardar e processar números. Para representar som
digitalmente, fazemos o seguinte: \textbf{medimos a pressão do ar (ou,
equivalentemente, a posição de um alto-falante) muitas vezes por segundo, e
guardamos cada medida como um número}.

\begin{newword}
\nova{Amostra} (\emph{sample}) Um único número que representa \enquote{qual
era a pressão do ar (ou posição do alto-falante) neste instante exato}. Uma
amostra sozinha não tem som nenhum --- é só um número, como uma única foto
de um vídeo.
\end{newword}

\begin{newword}
\nova{Taxa de amostragem} (\emph{sample rate}) Quantas amostras são medidas
(ou geradas) por segundo. Este projeto usa \textbf{48.000 amostras por
segundo} (48~kHz) --- o mesmo padrão usado em vídeo profissional e na
maioria dos equipamentos de áudio modernos. Isso significa: para produzir 1
segundo de som, o programa precisa calcular 48.000 números, em sequência,
um após o outro.
\end{newword}

Por que 48.000 e não, digamos, 100? Porque, para reconstituir uma onda com
fidelidade, é preciso amostrá-la a uma taxa \textbf{pelo menos duas vezes
maior} que a frequência mais alta que você quer representar (isso é um
resultado matemático chamado Teorema de Nyquist--Shannon). O ouvido humano
escuta até cerca de 20.000~Hz; 48.000 amostras por segundo dão folga
confortável acima disso.

Isso também nos dá um número muito importante para os próximos capítulos: a
cada amostra, o programa tem $\frac{1}{48000}$ de segundo --- cerca de
\textbf{20,8 microssegundos} ($0{,}0000208$ segundos) --- para calcular o
próximo número antes que ele precise ser enviado à placa de som. Repita
isso 48 mil vezes, e você tem um segundo de piano tocando. É essa corrida
contra o tempo, repetida sem parar, que faz o capítulo~15 (sobre o programa
\enquote{nunca poder travar}) ser tão importante.

% =============================================================================
\chapter{A grande sacada: uma corda é uma fila de espera}

\opening{C}{hegamos} à ideia central do projeto inteiro --- a peça de
engenharia que torna possível simular uma corda de piano em tempo real, sem
precisar de um supercomputador.

\section*{A analogia do eco no corredor}

Imagine um corredor comprido, com uma parede em cada ponta. Você grita uma
vez na entrada. O som viaja até a parede do fundo, ecoa, volta, ecoa de novo
na parede de onde você está, volta a ir\ldots{} e assim por diante, ficando
mais fraco a cada ida e volta, até sumir. O que você ouve, com o tempo, é
uma série de ecos regularmente espaçados, cada um mais fraco que o
anterior.

Uma corda de piano vibrando é \textbf{exatamente esse fenômeno}, só que
muito mais rápido: quando o martelo bate na corda, ele cria uma perturbação
que viaja de uma ponta da corda até a outra, reflete, volta, reflete de
novo --- e esse \enquote{ir e vir} que se repete centenas ou milhares de
vezes por segundo é o que você ouve como uma nota musical. A física formal
por trás disso (a \enquote{equação de onda}) mostra algo bonito: a vibração
de uma corda pode sempre ser descrita como \textbf{duas ondas viajando em
direções opostas}, quicando entre as duas pontas fixas.

\section*{Transformando isso em código: a linha de atraso}

Esse fato --- \enquote{é só uma onda indo e voltando} --- é ótimo porque
uma \enquote{onda indo e voltando} é exatamente o que um computador já sabe
representar naturalmente: uma \textbf{fila de espera com um número fixo de
posições}, onde a cada passo (a cada amostra) o número mais antigo sai por
um lado, todos os outros andam uma posição, e um número novo entra pelo
outro lado. Isso, em programação, se chama \textbf{linha de atraso}
(\emph{delay line}).

\begin{newword}
\nova{Linha de atraso} (\emph{delay line}) Uma fila de números de tamanho
fixo. A cada passo, um número novo entra, e o número mais antigo sai ---
como uma fila de pessoas esperando ônibus, onde uma pessoa nova entra na
fila atrás e a da frente sobe no ônibus a cada minuto.
\end{newword}

Em vez de calcular \enquote{qual é a altura de cada pontinho da corda, em
cada instante} (o que exigiria recalcular \emph{milhares} de pontos, 48 mil
vezes por segundo --- caro demais até para computadores modernos, se você
tem 88 teclas tocando ao mesmo tempo), o projeto guarda só \textbf{uma fila
circulando}, do tamanho certo, representando a onda viajando de ida e volta
pela corda. Isso é chamado de \textbf{guia de onda digital} (\emph{digital
waveguide} --- a técnica, publicada por Kevin Karplus e Alex Strong em
1983 e refinada depois por Julius O.\ Smith, é a base de todo
\texttt{piano\_core}, o coração deste projeto).

O custo computacional de processar essa fila é \textbf{o mesmo, não importa
quão grave ou aguda seja a corda} --- é isso que torna possível tocar as 88
teclas de um piano ao mesmo tempo em tempo real, mesmo num laptop comum.

\section*{Quantas posições tem a fila?}

Aqui é onde a frequência da nota (capítulo~3) entra na conta. Uma onda
completa \enquote{ida e volta} pela corda corresponde a exatamente um ciclo
da vibração. Então o tamanho da fila (quantas amostras ela precisa guardar)
é:

\[
\text{tamanho da fila} = \frac{\text{taxa de amostragem}}{\text{frequência da nota}}
\]

Para A4 (440~Hz), a 48.000 amostras por segundo:

\[
\text{tamanho da fila} = \frac{48\,000}{440} \approx 109{,}1 \text{ amostras}
\]

Ou seja: a fila que representa a corda de A4 tem pouco mais de 109
posições. Para a nota mais grave do piano, A0 (27,5~Hz), a fila tem
$\frac{48\,000}{27{,}5} \approx 1745$ posições --- uma fila bem mais longa,
porque uma corda grave é (física e metaforicamente) \enquote{mais comprida}
e o som demora mais para dar a volta completa.

Repare que o número não deu exatamente 109 --- deu 109,1. Essa fração é
importante, e o capítulo~9 explica por quê (e como o projeto lida com ela
sem desafinar a nota).

\section*{O resto do laço}

Uma fila sozinha, recirculando um número para sempre, tocaria a nota
\textbf{para sempre}, sem nunca morrer --- o que não é o que acontece numa
corda de verdade. Faltam ainda duas peças, que os próximos dois capítulos
explicam: algo que \textbf{tire energia} da fila a cada volta (o som
precisa morrer com o tempo, capítulo~8), e algo que \textbf{inicie} a
vibração de um jeito parecido com uma martelada de feltro, não um empurrão
genérico (capítulo~11).

% =============================================================================
\chapter{Por que a nota morre: o filtro que come agudos}

\opening{N}{uma} corda ideal e sem atrito, a energia da martelada ficaria
circulando na fila para sempre --- a nota tocaria infinitamente, o que
obviamente não é o que acontece num piano de verdade. Na vida real, a corda
perde energia a cada ida e volta: um pouco escapa pela ponte para a caixa
de ressonância (de propósito --- é assim que você ouve o som!), e um pouco
se perde em atrito interno do próprio metal e do ar ao redor.

O detalhe físico interessante --- e crucial para soar como um piano de
verdade, não como um bipe genérico --- é que \textbf{essa perda de energia
não é igual para todos os harmônicos}. Os harmônicos mais agudos (as
vibrações mais rápidas, capítulo~5) perdem energia muito mais depressa que
os graves. É por isso que uma nota de piano começa com um ataque
brilhante, cheio de agudos, e vai ficando cada vez mais \enquote{redonda} e
grave conforme soa --- os agudos vão embora primeiro, sobra só o
fundamental por último.

\begin{newword}
\nova{Filtro passa-baixa} (\emph{lowpass filter}) Um mecanismo que deixa
frequências baixas passarem quase sem mudança, mas reduz (atenua)
frequências altas. É exatamente o oposto de um cobertor grosso jogado sobre
uma caixa de som: os graves (o \enquote{bum bum} do baixo) você ainda ouve
do lado de fora, os agudos (o \enquote{tss tss} do prato de bateria) somem
quase por completo --- o cobertor está agindo como um filtro passa-baixa
físico.
\end{newword}

A solução do projeto: dentro da fila de espera (a linha de atraso do
capítulo~7), antes do número voltar a circular, ele passa por um filtro
passa-baixa simples --- implementado como \texttt{filter::OnePoleLowpass}
no código --- que reduz um pouquinho a \enquote{agudeza} do número a cada
volta. Feito isso milhares de vezes por segundo, o resultado acumulado é
exatamente o comportamento real: brilhante no início, cada vez mais opaco,
até silêncio.

\begin{newword}
\nova{\enquote{Um polo} (\emph{one-pole})} A versão mais simples possível
de um filtro --- precisa de só uma multiplicação e uma soma por amostra
para funcionar. \enquote{Polo} é um termo técnico de processamento de
sinais; o que importa saber é que, apesar de simples, um filtro de um polo
já é suficiente para capturar o comportamento essencial de \enquote{agudos
morrem mais rápido que graves} que uma corda de verdade exibe.
\end{newword}

O tanto de energia que se perde a cada volta é controlado, no código, por
um parâmetro chamado \texttt{damping}.

\begin{newword}
\nova{Amortecimento} (\emph{damping}) O quanto de energia uma corda perde
por unidade de tempo; mais amortecimento = a nota morre mais rápido.
\end{newword}

É esse número --- junto de um segundo parâmetro, \texttt{sustain} --- que
faz uma corda grave soar por 30 a 40 segundos e uma corda aguda por apenas
1 a 2 segundos, exatamente como num piano de verdade.

% =============================================================================
\chapter{O ajuste fino: por que o laço é um pouquinho mais curto}

\opening{V}{oltando} ao capítulo~7: calculamos que a fila para A4 devia ter
cerca de $109{,}1$ posições. Só que uma fila de espera só pode ter um
número \textbf{inteiro} de posições --- não existem \enquote{0,1 de uma
posição}. Se simplesmente arredondássemos para 109, a nota tocaria
ligeiramente mais aguda que 440~Hz; se arredondássemos para cima,
ligeiramente mais grave. Qualquer um dos dois desafinaria o piano --- pouco
nas notas graves, perceptivelmente nas agudas (quanto mais aguda a nota,
menor é a fila, e um erro de \enquote{1 posição inteira} pesa
proporcionalmente mais).

A solução tecnicamente correta é permitir uma fila de tamanho
\textbf{fracionário} --- por exemplo, calcular a posição \enquote{109,1}
como uma média ponderada entre a posição 109 e a posição 110. Esse
pedacinho de matemática (chamado \textbf{interpolação}) é o que garante que
qualquer nota, não só as que \enquote{dão números redondos} de amostras,
saia afinada.

Tem ainda uma segunda sutileza. O filtro passa-baixa do capítulo~8 não é
instantâneo --- ele próprio introduz um pequeno atraso extra no sinal
(assim como esperar um cobertor \enquote{amortecer} um som leva uma fração
de segundo, não é instantâneo). Se esse atraso extra não fosse descontado,
o laço inteiro ficaria um pouquinho mais longo do que deveria, e a nota
sairia ligeiramente mais grave do que a frequência pedida. A correção é
simples de enunciar, ainda que a matemática exata do \enquote{quanto} venha
de como o filtro é construído:

\[
\text{tamanho real da fila} = \text{período da nota} - \text{atraso do filtro} - 1
\]

Esse ajuste está implementado em \texttt{PluckedString::new}, e o projeto
tem um teste automatizado
(\texttt{loop\_delay\_is\_close\_to\_the\_period}) que verifica, toda vez
que o código muda, que a nota A4 realmente sai a menos de
\textbf{1,5~cents} de 440~Hz.

\begin{newword}
\nova{Cents} A unidade que músicos usam para medir desafinação bem fina.
Um semitom (o menor \enquote{degrau} do piano, capítulo~3) vale 100 cents.
1,5 cents é uma fração tão pequena de um semitom (1,5\,\%) que praticamente
nenhum ouvido humano --- nem afinadores profissionais, na maioria dos casos
--- consegue perceber a diferença. É esse o padrão de precisão que o
projeto garante, e verifica automaticamente, para cada nota.
\end{newword}

% =============================================================================
\chapter{Cordas de verdade são rígidas: a inarmonicidade}

\opening{N}{o} capítulo~5 dissemos que os harmônicos de uma corda ficam em
múltiplos exatos da fundamental: $2\times$, $3\times$, $4\times$ etc. Isso
é verdade para uma corda \textbf{ideal} --- perfeitamente flexível, sem
nenhuma resistência a dobrar. Cordas de piano de verdade, porém, são feitas
de \textbf{aço grosso e razoavelmente rígido} (principalmente nas notas
graves), e essa rigidez faz os harmônicos ficarem \textbf{ligeiramente mais
agudos} do que os múltiplos exatos previstos --- um efeito chamado
\textbf{inarmonicidade}, descrito matematicamente pela primeira vez pelo
físico Harvey Fletcher em 1964.

\begin{newword}
\nova{Inarmonicidade} O quanto os harmônicos de um instrumento real se
desviam de serem múltiplos exatos da fundamental, por causa da rigidez
física da corda. É pequena, mas real --- e é uma das razões pelas quais um
piano de verdade soa \enquote{vivo} e um bipe eletrônico perfeitamente
harmônico soa \enquote{morto} ou \enquote{artificial}.
\end{newword}

A fórmula de Fletcher para onde cada harmônico realmente cai é:

\[
f_n \approx n \, f_1 \sqrt{1 + B n^2}
\]

Não é preciso decorar a fórmula --- o que ela diz, em palavras, é: \enquote{o
harmônico número $n$ fica um pouquinho mais agudo do que $n$ vezes a
fundamental, e esse desvio cresce com o quadrado de $n$ (então harmônicos
mais altos desviam proporcionalmente mais), multiplicado por um número $B$
que descreve o quão rígida é aquela corda específica}. Cordas grossas e
curtas (as graves) têm um $B$ maior; cordas finas e longas (as agudas) têm
um $B$ menor.

Este é, aliás, o motivo pelo qual afinadores de piano profissionais não
afinam o instrumento em oitavas \emph{matematicamente} exatas --- eles
\enquote{esticam} ligeiramente as oitavas graves e agudas para compensar a
inarmonicidade e fazer o piano soar afinado \emph{aos ouvidos}, não apenas
no papel. Esse fenômeno, bem documentado na literatura de acústica
musical, é a mesma razão física que este projeto reproduz.

\subsection*{Como isso vira código}

Reproduzir esse efeito dentro da \enquote{fila de espera} do capítulo~7 usa
uma peça chamada \textbf{cascata de filtros allpass}
(\texttt{piano\_core::dispersion::DispersionCascade}), publicada por David
Jaffe e Julius O.\ Smith em 1983.

\begin{newword}
\nova{Filtro allpass} Ao contrário do filtro passa-baixa do capítulo~8 (que
muda o \emph{volume} de diferentes frequências), um filtro allpass deixa o
volume de \textbf{todas} as frequências exatamente igual --- mas atrasa
cada frequência por um tempo ligeiramente diferente. Isso é exatamente o
que \enquote{esticar} os harmônicos para fora da posição exata significa:
cada harmônico chega um pouquinho \enquote{fora de sincronia} com onde
estaria numa corda ideal, sem que nenhum fique mais alto ou baixo em
volume.
\end{newword}

Encadear vários desses filtros (\enquote{uma cascata}) dentro do laço da
corda aproxima, com boa precisão, a curva de inarmonicidade real. Quantos
filtros são necessários varia por registro --- cordas graves precisam de
mais seções (cerca de 8 para a nota mais grave do piano) e cordas agudas
precisam de quase nenhuma (0 a 1 para as notas mais agudas), porque o
efeito de inarmonicidade é fisicamente mais forte nas cordas grossas e
curtas do que nas finas e longas. Usar sempre o número máximo de seções,
mesmo onde o ouvido não notaria diferença, gastaria processador à toa ---
por isso o projeto escala esse número por registro (documentado como
decisão de performance \texttt{PERF-005}).

% =============================================================================
\chapter{O martelo de feltro: como uma martelada vira som}

\opening{A}{té} agora falamos do que acontece \textbf{depois} que a corda
já está vibrando. Falta a pergunta mais básica: \textbf{como a vibração
começa?}

A técnica clássica de Karplus e Strong (1983), da qual este projeto parte,
usa uma solução simples: enche a fila de espera inteira com \textbf{ruído
aleatório}.

\begin{newword}
\nova{Ruído} Um sinal que contém, em proporções parecidas, todas as
frequências ao mesmo tempo --- como a \enquote{estática} de uma TV fora de
sintonia; matematicamente é uma forma prática de \enquote{acender} todos os
harmônicos de uma corda de uma vez.
\end{newword}

Fisicamente, isso corresponde a \enquote{empurrar a corda inteira para uma
posição aleatória de uma vez} --- o que é uma boa aproximação para um
instrumento \textbf{beliscado}, como um violão. Mas um piano não é
beliscado, é \textbf{golpeado por um martelo de feltro}, e é exatamente
esse detalhe que faz o piano soar como piano.

\subsection*{Por que a força da martelada muda o timbre, não só o volume}

Toque uma tecla de piano bem de leve, depois toque a mesma tecla com toda
a força. Você não vai ouvir \enquote{o mesmo som, só que mais alto} --- vai
ouvir um som mais \textbf{brilhante}, com mais agudos, quando toca forte.
Esse é um dos traços mais característicos de um piano de verdade, e a razão
física é o comportamento do próprio feltro do martelo.

O feltro que cobre o martelo funciona como uma mola, mas uma \textbf{mola
não-linear}: quanto mais você a comprime, mais dura ela fica (diferente de
uma mola comum, cuja dureza não muda). Esse comportamento foi medido e
descrito por Antoine Chaigne e Anders Askenfelt em 1994, com a fórmula:

\[
F = K \, x^{p} \qquad \text{(com } p \text{ entre 2 e 3, aproximadamente)}
\]

Em palavras: a força $F$ que o martelo exerce sobre a corda cresce
\textbf{mais rápido que proporcionalmente} conforme a compressão $x$ do
feltro aumenta (por isso o expoente $p$ é maior que 1 --- se fosse $p=1$,
dobrar a compressão dobraria a força; com $p$ entre 2 e 3, dobrar a
compressão multiplica a força por 4 a 8 vezes). Uma martelada forte
comprime mais o feltro, o que o deixa efetivamente mais \enquote{duro}
durante aquele contato específico --- um feltro mais duro transmite um
golpe \textbf{mais curto e mais brusco}, e um golpe mais brusco excita
harmônicos mais agudos com mais força. Uma martelada fraca comprime pouco,
o feltro fica \enquote{macio}, o golpe é mais longo e suave, e o som sai
mais quente, com menos agudos.

Este projeto implementa esse comportamento em
\texttt{piano\_core::hammer::simulate\_contact}: dada a velocidade com que
a tecla foi pressionada, o código simula o contato do martelo com a mola
não-linear de feltro e produz um \textbf{envelope de força} (a forma da
martelada ao longo do tempo, que dura menos de 5 milésimos de segundo).
Esse envelope é usado para \emph{moldar} o ruído inicial que preenche a
fila de espera --- em vez de simplesmente \enquote{mais forte = mais
alto}, a forma do ataque muda com a velocidade.

\begin{newword}
\nova{Honestidade sobre o que ainda falta aqui} O modelo atual simula o
lado do \emph{martelo} dessa interação --- como o feltro se comprime e
empurra --- mas, durante os poucos milissegundos de contato, ainda não
simula a corda \enquote{empurrando de volta} o martelo em tempo real (o
problema físico completo, chamado de contato acoplado). Isso é uma
simplificação documentada de propósito (rastreada como \texttt{PERF-007}
na documentação técnica), e é exatamente o item no topo da lista de
melhorias futuras do capítulo~17.
\end{newword}

% =============================================================================
\chapter{Mais de uma corda por tecla: uníssono e batimento}

\opening{A}{qui} vai um fato que surpreende muita gente: \textbf{a maioria
das teclas de um piano não tem apenas uma corda} --- tem duas ou três,
todas afinadas quase (mas não exatamente) na mesma frequência, e um único
martelo bate nas cordas de uma tecla ao mesmo tempo. Um piano moderno
típico tem:

\begin{itemize}
\item \textbf{1 corda} por tecla nas notas mais graves (mais grossas, mais
caras, mais espaço ocupado --- não cabem duas);
\item \textbf{2 cordas} por tecla na região média-grave;
\item \textbf{3 cordas} por tecla em toda a região média-aguda e aguda.
\end{itemize}

Este projeto usa essa mesma convenção --- 12 teclas com 1 corda, 18 com 2 e
58 com 3 --- o que soma \textbf{222 cordas efetivas} sendo processadas ao
vivo, mesmo com só 88 teclas.

\begin{newword}
\nova{Uníssono} (\emph{unison}) O grupo de 1 a 3 cordas que uma única
tecla controla. As cordas de um mesmo uníssono são afinadas \emph{quase}
iguais, mas não perfeitamente --- normalmente há uma diferença de poucos
\enquote{cents} (capítulo~9) entre elas, de propósito.
\end{newword}

\subsection*{Por que desafinar de propósito?}

Se as duas ou três cordas de uma tecla fossem afinadas \textbf{exatamente}
iguais, elas vibrariam perfeitamente em sincronia, e o resultado soaria
idêntico a uma única corda --- apenas mais alto. Mas quando estão afinadas
com uma diferença bem pequena, acontece um fenômeno físico chamado
\textbf{batimento}:

\begin{newword}
\nova{Batimento} (\emph{beating}) Quando duas ondas sonoras de frequências
muito próximas (mas não idênticas) tocam juntas, elas alternadamente se
reforçam e se cancelam, produzindo um efeito de \enquote{pulsar} no volume
--- um \enquote{uauauauau} lento, cuja velocidade é exatamente igual à
diferença entre as duas frequências. Você pode ouvir isso muito
claramente quando duas cordas de violão quase afinadas tocam juntas: em
vez de um som estável, ouve-se um \enquote{bater} rítmico que vai sumindo
conforme se ajusta a afinação.
\end{newword}

Esse \enquote{bater} é parte do que dá a um piano de verdade seu caráter
rico e \enquote{vivo}, em vez de um som mais \enquote{seco} e sintético.

\subsection*{O envelope de dois estágios}

Tem mais uma consequência física interessante, medida e descrita
formalmente pelo físico Gabriel Weinreich em 1977: as cordas de um mesmo
uníssono não vibram de forma independente --- elas estão mecanicamente
conectadas através da ponte (capítulo~4), que embora seja rígida, cede um
pouquinho. Essa conexão parcial faz o som de uma nota de piano decair em
\textbf{dois estágios}, não um só:

\begin{enumerate}[leftmargin=1.6em]
\item \textbf{Pré-decaimento} (rápido): logo após a martelada, as cordas
ligeiramente desafinadas batem entre si e perdem energia relativa umas às
outras rapidamente --- essa fase soa mais \enquote{cheia} e um pouco
instável.
\item \textbf{Cauda} (mais lenta): depois que essa energia diferencial se
dissipa, as cordas se estabilizam num modo compartilhado que decai numa
taxa parecida com a de uma única corda \enquote{natural} --- essa é a fase
mais longa e estável, o \enquote{sustain} que você ouve segurando a tecla.
\end{enumerate}

Este projeto implementa exatamente esse mecanismo (não um envelope
artificial \enquote{desenhado à mão} para parecer parecido) em duas
camadas:

\begin{itemize}
\item \textbf{Acoplamento local} (\texttt{piano\_core::unison}): as 1 a 3
cordas da mesma tecla se misturam entre si a cada amostra individual ---
é barato de calcular porque essas cordas já são processadas juntas de
qualquer forma.
\item \textbf{Acoplamento global} (\texttt{piano\_core::bridge::BridgeBus},
capítulo~13): a interação entre teclas \emph{diferentes}, através da
ponte compartilhada de todo o piano.
\end{itemize}

O projeto tem inclusive um teste automatizado que \textbf{mede} esse
efeito (não apenas confia que o código \enquote{deveria} produzir isso):
ele renderiza uma nota de três cordas, mede a velocidade do decaimento logo
após o ataque comparada com a velocidade depois que o batimento já se
estabilizou, e confirma que a primeira é visivelmente mais rápida que a
segunda --- a assinatura exata que o modelo de Weinreich prevê. Esse teste,
durante o desenvolvimento, pegou dois bugs reais de implementação antes que
chegassem ao código final --- um exemplo de como medir o comportamento
físico é mais confiável do que apenas ler o código e \enquote{achar} que
está certo.

% =============================================================================
\chapter{O pedal mágico: ressonância simpática}

\opening{S}{ente-se} ao piano, pise no pedal direito (sustentação) sem
tocar nenhuma tecla, e cante ou grite uma nota perto das cordas. Você vai
ouvir o piano \enquote{responder} --- alguma corda dentro dele começa a
vibrar sozinha, mesmo sem ninguém tê-la tocado. Isso é \textbf{ressonância
simpática}, e é um dos efeitos mais mágicos (e mais difíceis de simular
bem) de um piano de verdade.

\begin{newword}
\nova{Ressonância simpática} Quando um objeto vibrante (uma corda, uma
voz) faz \emph{outro} objeto próximo, capaz de vibrar na mesma frequência
ou numa relacionada, começar a vibrar também --- sem nenhum contato
direto, só através do ar ou de uma estrutura compartilhada. É o mesmo
princípio por trás de empurrar um balanço no ritmo certo para fazê-lo ir
cada vez mais alto: pequenos empurrões, no momento certo, se acumulam.
\end{newword}

Num piano de verdade, isso acontece porque \textbf{todas as cordas do
instrumento} estão fisicamente conectadas através da mesma ponte e da
mesma caixa de ressonância. Normalmente, os abafadores de feltro mantêm as
cordas que você não está tocando \enquote{mudas} --- mas quando você pisa
no pedal direito, \textbf{todos os abafadores levantam ao mesmo tempo}, e
qualquer corda cuja frequência (ou harmônico) combine com o que está soando
pode \enquote{pegar carona} na vibração e começar a soar também, mesmo sem
ter sido golpeada.

\subsection*{Como reproduzir isso sem explodir o custo computacional}

A forma fisicamente mais completa de simular isso seria calcular como
\textbf{cada uma das 222 cordas efetivas afeta todas as outras 221},
individualmente --- o que dá mais de 49 mil pares de interação,
recalculados a cada amostra, 48 mil vezes por segundo. Isso é
computacionalmente inviável em tempo real, mesmo em computadores potentes.

A solução do projeto (\texttt{piano\_core::bridge::BridgeBus}) é mais
elegante: todas as cordas escrevem sua vibração num único
\textbf{barramento compartilhado} --- pense nele como um \enquote{balde}
onde cada corda despeja um pouco do seu som, e de onde cada corda também lê
uma média de tudo que está no balde. Isso captura o essencial do efeito
real (cordas \enquote{sentem} o que as outras estão fazendo) com um custo
que cresce de forma administrável --- em vez de crescer explosivamente com
o número de cordas.

\begin{newword}
\nova{Simplificação, dita com honestidade} O barramento único usa o
\textbf{mesmo} ganho de acoplamento para todas as cordas, e não modela como
a ponte de um piano real transmite som de forma diferente dependendo de
onde, na madeira, cada corda está apoiada --- isso exigiria medições de um
instrumento físico específico, algo que o projeto delibera não fazer sem
antes decidir explicitamente que fontes de dado são aceitáveis
(capítulo~17). É uma aproximação de engenharia, declarada abertamente na
documentação técnica, não apresentada como mais fiel do que realmente é.
\end{newword}

% =============================================================================
\chapter{A caixa de ressonância: de onde vem o volume}

\opening{C}{omo} mencionado no capítulo~4, uma corda de piano sozinha é
fininha demais para empurrar ar suficiente e ser ouvida com volume real ---
quase todo o som que você de fato escuta vem da \textbf{caixa de
ressonância}, o grande painel de madeira que a ponte transmite a vibração
para.

A forma mais \enquote{fiel à física} de simular uma caixa de ressonância
seria gravar, de um piano real, como ele responde a um impulso curtíssimo
(um \enquote{clique}), e depois combinar matematicamente
(\enquote{convoluir}) essa gravação com o som de cada corda. \textbf{Este
projeto não pode fazer isso}, e não por falta de capacidade técnica: a
regra fundamental do projeto proíbe qualquer gravação de áudio real, para
sempre, em qualquer circunstância --- todo o som precisa nascer de
cálculo, nunca de amostras gravadas.

A alternativa escolhida é a \textbf{síntese modal}: representar a caixa de
ressonância como um conjunto de \enquote{modos} --- frequências
específicas nas quais a madeira naturalmente prefere vibrar, cada uma com
seu próprio tempo de decaimento e volume relativo, baseados em valores
típicos descritos na literatura acadêmica de acústica (não medidos de
nenhum piano específico).

\begin{newword}
\nova{Modo de ressonância} Assim como uma corda tem harmônicos
(capítulo~5), uma placa de madeira grande também tem frequências
preferenciais nas quais ela \enquote{gosta} de vibrar quando excitada ---
mas, ao contrário de uma corda (onde os harmônicos formam uma série
organizada de múltiplos), os modos de uma placa bidimensional são mais
irregulares e dependem do formato e material específicos da placa.
\end{newword}

O projeto usa 8 desses modos (\texttt{piano\_core::soundboard::Soundboard}),
somados ao sinal direto de cada corda (não substituindo-o) antes de virar
som final --- porque substituir completamente mudaria a afinação e o
brilho medidos e testados em capítulos anteriores, uma troca
(\emph{trade-off}) que o projeto documenta abertamente em vez de mudar
silenciosamente.

% =============================================================================
\chapter{Por que o programa nunca pode travar}

\opening{L}{embra} do capítulo~6? A cada $\frac{1}{48000}$ de segundo
(cerca de 20,8 microssegundos), o programa precisa ter calculado o próximo
número de áudio, para \textbf{todas} as vozes tocando simultaneamente, e
entregá-lo à placa de som. Se ele demorar mais que isso mesmo uma única
vez, o resultado não é \enquote{um pouquinho de atraso} --- é um
\textbf{estalo, clique ou silêncio audível} na música, porque a placa de
som não tem o que tocar naquele instante exato.

Isso impõe uma regra muito mais rígida do que a maioria dos programas de
computador precisa seguir. A parte do código que calcula o áudio (chamada
de \textbf{thread de tempo real}, porque roda numa linha de execução
separada e dedicada exclusivamente a isso) segue quatro regras absolutas,
verificadas automaticamente a cada mudança no código:

\begin{enumerate}[leftmargin=1.6em]
\item \textbf{Nunca aloca memória nova.} Pedir memória nova ao sistema
operacional pode, ocasionalmente, demorar um tempo imprevisível ---
inaceitável dentro da janela de 20,8 microssegundos. Toda a memória que
uma nota vai precisar é reservada com antecedência, na hora em que o
programa começa a rodar, não quando você aperta a tecla.
\item \textbf{Nunca espera por outra parte do programa} (nenhum
\enquote{cadeado} ou trava de sincronização). Se essa parte do código
tivesse que esperar outra parte terminar algo, e essa outra parte
estivesse, por qualquer motivo, atrasada, o áudio inteiro travaria com
ela.
\item \textbf{Nunca pode gerar um erro fatal.} Um programa comum, ao
encontrar uma situação inesperada, pode simplesmente parar com uma
mensagem de erro. Isso é inaceitável no meio de uma peça musical --- então
todo pedaço de código nessa parte do projeto é escrito de um jeito que
\textbf{sempre} devolve algum resultado válido, mesmo em casos extremos
(um número \enquote{infinito}, um número inválido, zero etc.).
\item \textbf{Nunca entra num laço sem fim garantido.} Todo cálculo tem um
número máximo de passos conhecido de antemão --- nada de \enquote{repita
até algo acontecer} sem um limite garantido.
\end{enumerate}

O projeto até usa uma ferramenta do próprio compilador da linguagem Rust
(chamada \texttt{no\_std}, e uma configuração que \textbf{proíbe} o uso das
funções \texttt{unwrap}, \texttt{expect} e \texttt{panic!} fora de testes)
para tornar \textbf{impossível}, não apenas malvisto, escrever código que
quebre essas regras sem querer --- o próprio compilador recusa compilar o
programa se alguém tentar. Isso é o que separa \enquote{um programa que
geralmente não trava} de \enquote{um programa que matematicamente não pode
travar por esses motivos}.

% =============================================================================
\chapter{Como as peças do projeto se encaixam}

\opening{T}{odo} o código está organizado em módulos independentes
(chamados, na linguagem Rust, de \textbf{crates}), cada um com uma única
responsabilidade clara, para que uma mudança numa parte não possa quebrar
outra por acidente:

\begin{center}
\begin{tcolorbox}[colback=white,colframe=rulecol,boxrule=0.6pt,arc=0pt,
  left=10pt,right=10pt,top=8pt,bottom=8pt,width=0.92\linewidth]
\ttfamily\small
\begin{verbatim}
piano-core     A fisica pura. Nenhuma peca deste modulo sabe se
               esta rodando num computador, num navegador ou
               num teste -- ele so recebe numeros e devolve numeros.
                 |
                 +-- piano-params   Nomes de notas, numeros MIDI, afinacao.
                 |     |
                 |     +-- piano-render   Renderiza para arquivos .wav
                 |     |     |
                 |     |     +-- piano-cli   O programa "piano"
                 |     |
                 |     +-- piano-audio    Toca ao vivo pelos alto-falantes
                 |     +-- piano-midi     Recebe teclados MIDI de verdade
                 |
                 +-- piano-wasm     Versao que roda dentro do navegador
\end{verbatim}
\end{tcolorbox}
\end{center}

A ideia central --- repetida por todo o projeto --- é: \textbf{o motor de
física (\texttt{piano-core}) não sabe nada sobre o mundo exterior}. Ele não
sabe se está sendo tocado por um teclado de computador, por um piano MIDI
de verdade conectado por USB, ou por uma página de navegador. Só as
camadas \emph{fora} dele (\texttt{piano-audio}, \texttt{piano-midi},
\texttt{piano-wasm}, \texttt{piano-cli}) sabem lidar com o \enquote{mundo
real} --- placas de som, cabos USB, páginas web. Isso é o que permite o
\textbf{mesmo} código de física rodar de forma idêntica em três lugares
completamente diferentes, sem nenhuma parte especial para cada um.

% =============================================================================
\chapter{O que ainda falta para chegar perto de um piano de verdade}

\opening{E}{ste} projeto não é um protótipo incompleto --- cada mecanismo
descrito nos capítulos~7 a~14 já está implementado, testado e verificado
(não apenas planejado). Mas um piano de verdade tem ainda mais detalhes do
que os já cobertos, e este capítulo lista os que faltam, honestamente, do
maior impacto perceptível para o menor:

\begin{table}[h!]
\centering
\small
\begin{tabularx}{\linewidth}{p{3.6cm}Xl}
\toprule
\textbf{O que falta} & \textbf{Por que importa} & \textbf{Marco} \\
\midrule
Contato martelo-corda totalmente acoplado &
Hoje o martelo (capítulo~11) empurra a corda, mas a corda ainda não
\enquote{empurra de volta} o martelo durante a martelada --- a versão
completa muda sutilmente o ataque de cada nota & M9 \\
\addlinespace
Posição da martelada e curvas por tecla &
Onde exatamente o martelo bate ao longo da corda muda o timbre; hoje isso
ainda não varia por nota & M9 \\
\addlinespace
Pedal una corda (esquerdo) e sostenuto (do meio) &
Hoje só existe o pedal de sustentação (direito); um piano de verdade tem
três pedais, cada um com um efeito diferente & M10 \\
\addlinespace
Ressonância simpática mais rica &
O barramento único (capítulo~13) já funciona, mas uma versão mais fiel
diferenciaria o quanto cada corda \enquote{sente} as outras & M11 \\
\addlinespace
Caixa de ressonância mais detalhada &
8 modos (capítulo~14) já dão corpo ao som; um banco maior aproximaria
ainda mais a complexidade real da madeira & M11 \\
\addlinespace
Ruídos mecânicos &
Um piano de verdade não é só corda --- o mecanismo em si faz ruído, sempre
sintetizado, nunca gravado & M12 \\
\addlinespace
Modelar um piano específico de verdade &
Hoje o projeto usa números \enquote{típicos} da literatura, não medidos de
um instrumento físico específico; exigiria decidir, antes de qualquer
código, que dado medido é aceitável & M13 \\
\addlinespace
Um plugin de estúdio (CLAP) &
Permitiria usar o piano dentro de programas como Ableton, Logic ou Reaper
& M14 \\
\bottomrule
\end{tabularx}
\end{table}

Esses marcos (M9 a M14) já existem como \textbf{issues públicas no GitHub}
do projeto, para quem quiser acompanhar ou contribuir --- não são apenas
ideias soltas, são trabalho planejado e rastreável.

Vale reforçar: mesmo com toda essa lista de \enquote{o que falta}, o piano
\textbf{já funciona de verdade hoje} --- os capítulos~7 a~14 não são
teoria, são o comportamento real do programa que você pode compilar e
tocar agora mesmo (veja \texttt{docs/pt-BR/LEIA-ME.md} para o passo a
passo de instalação e uso).

% =============================================================================
\backmatter
\pagestyle{fancy}

\chapter*{Glossário completo}
\addcontentsline{toc}{chapter}{Glossário completo}
\markboth{Glossário completo}{}

\noindent Todas as palavras novas deste documento, num só lugar, para
consulta rápida.

\begin{center}
\small
\begin{tabularx}{\linewidth}{p{4.4cm}X}
\toprule
\textbf{Termo} & \textbf{Significado} \\
\midrule
Onda & Uma perturbação que se propaga por um meio sem que o meio em si viaje junto. \\
Frequência & Quantas vezes por segundo algo se repete, medida em hertz (Hz). \\
Hertz (Hz) & Unidade de frequência: \enquote{vezes por segundo}. \\
Nota musical & Um som cuja frequência é constante e bem definida durante o tempo em que soa. \\
Oitava & O intervalo entre uma frequência e o dobro dela. \\
Semitom & O menor degrau entre notas num piano; 12 semitons formam uma oitava. \\
A4 / Lá4 & Nota de referência do piano, fixada em 440~Hz. \\
Harmônico / parcial / overtone & Cada uma das vibrações simultâneas de uma corda, em múltiplos da fundamental. \\
Fundamental & O primeiro harmônico --- a corda vibrando inteira, de ponta a ponta. \\
Timbre & A \enquote{cor} característica de um instrumento, determinada pela mistura de harmônicos. \\
Amostra (\emph{sample}) & Um único número medindo a pressão do ar (ou posição do alto-falante) num instante. \\
Taxa de amostragem & Quantas amostras são geradas por segundo (48.000 neste projeto). \\
Linha de atraso (\emph{delay line}) & Uma fila de números de tamanho fixo, onde um número novo entra e o mais antigo sai a cada passo. \\
Guia de onda digital & A técnica de representar uma corda vibrando como uma linha de atraso circulando. \\
Filtro passa-baixa & Um mecanismo que reduz frequências altas e deixa as baixas passarem quase sem mudança. \\
Amortecimento (\emph{damping}) & O quanto de energia uma corda perde por unidade de tempo. \\
Cents & Unidade fina de afinação; 100 cents formam um semitom. \\
Inarmonicidade & O desvio dos harmônicos de uma corda real em relação a múltiplos exatos da fundamental, causado pela rigidez física da corda. \\
Filtro allpass & Um filtro que atrasa frequências diferentes por tempos diferentes, sem mudar o volume de nenhuma. \\
Ruído & Um sinal que contém várias frequências ao mesmo tempo, em proporções parecidas. \\
Uníssono (\emph{unison}) & O grupo de 1 a 3 cordas que uma mesma tecla de piano controla. \\
Batimento (\emph{beating}) & O \enquote{pulsar} de volume que surge quando dois sons de frequências muito próximas tocam juntos. \\
Ressonância simpática & Quando um objeto vibrante faz outro, próximo, começar a vibrar também, sem contato direto. \\
Modo de ressonância & Uma frequência preferencial na qual uma estrutura naturalmente vibra. \\
Ponte (\emph{bridge}) & A peça que transmite a vibração de uma corda para a caixa de ressonância. \\
Modelagem física & Descrever um objeto real como equações e resolvê-las em tempo real para produzir seu comportamento. \\
Sampler & Um instrumento que reproduz gravações prontas em vez de calcular o som. \\
\bottomrule
\end{tabularx}
\end{center}

\chapter*{Para saber mais}
\addcontentsline{toc}{chapter}{Para saber mais}
\markboth{Para saber mais}{}

\opening{E}{ste} documento é a versão didática e sem pré-requisitos de
material técnico mais denso que já existe no projeto, em inglês, para quem
quiser ir mais fundo:

\begin{itemize}
\item \textbf{\texttt{docs/PHYSICS.md}} --- o mapeamento completo, com mais
detalhes técnicos e as referências acadêmicas exatas por trás de cada
mecanismo descrito neste documento.
\item \textbf{\texttt{docs/ARCHITECTURE.md}} --- como o código está
organizado por dentro.
\item \textbf{\texttt{docs/PERFORMANCE.md}} --- as medições reais de
desempenho por trás de cada decisão de engenharia mencionada aqui.
\item \textbf{\texttt{docs/ROADMAP.md}} --- todos os marcos do projeto,
passados e futuros.
\item \textbf{\texttt{docs/pt-BR/LEIA-ME.md}} --- o guia de instalação e
uso em português, passo a passo, para quem quer simplesmente tocar o piano
agora.
\item Os artigos científicos originais citados ao longo deste documento
(Karplus e Strong 1983; Jaffe e Smith 1983; Fletcher 1964; Chaigne e
Askenfelt 1994; Weinreich 1977) estão listados, com título completo e onde
encontrá-los, em \texttt{docs/PRIOR-ART.md}.
\end{itemize}

Se depois de ler tudo isso uma pergunta ainda ficou sem resposta, ela
provavelmente merece virar uma \emph{issue} no GitHub do projeto --- é
assim que esta documentação, e o próprio piano, continuam melhorando.

\end{document}
